\subsection{Effective Expansion Rate in Voids}
  \label{subsec:void_effect}

  Cosmic voids represent the deepest and most extended low-density regions in the
  late Universe.
  Within voids, baryonic and total matter densities are strongly suppressed,
  and typical gravitational accelerations fall below the saturation scale
  $a_\star$ over large volumes~\cite{Keenan2013KBC}.

  In this regime, the effective description operates close to the limited descriptive resolution in low-density
  environments.
  The relation between local kinematics and the underlying background expansion
  is therefore modified at the level of inference rather than dynamics.
  To leading order, this effect can be captured by an environment-dependent
  renormalization of the inferred expansion rate,
  \begin{equation}
    H_{\mathrm{loc}}
    =
    H_{\mathrm{global}}
    \left[
      1
      +
      \delta_{\mathrm{eff}}
    \right] ,
    \quad \text{with} \quad
    \delta_{\mathrm{eff}} = \delta_\rho \cdot f\!\left(
                                                   \frac{g_{\mathrm{void}}}{a_\star}
    \right),
    \label{eq:hloc}
  \end{equation}
  where
  \(\delta_\rho = 1 - \rho_{\mathrm{void}}/\rho_{\mathrm{global}}\)
  is the density contrast of the void,
  and
  \(g_{\mathrm{void}} = G M_{\mathrm{void}} / R_{\mathrm{void}}^2\)
  is the characteristic gravitational field strength within the void.
  The function \(f(x)\) is a dimensionless saturation function with the properties:
  \begin{itemize}
    \item \(f(x) \rightarrow 0\) for \(x \gg 1\) (unsaturated regime),
    \item \(0 < f(x) < f_{\mathrm{max}}\) for \(x \ll 1\) (saturated regime),
    \item \(f_{\mathrm{max}} \lesssim 0.1\), reflecting the maximal projection deficit in deep voids
    (Section 4.3 of~\cite{Beau2026c}).
  \end{itemize}
  The saturation scale \(a_\star\) is fixed independently by galaxy rotation curve data
  (Section~\ref{subsec:universality}).

  For typical voids (\(\delta_\rho \approx 0.3\), \(g_{\mathrm{void}} \ll a_\star\)),
  Eq.~\eqref{eq:hloc} yields \(\Delta H/H \sim 0.03\)--\(0.08\),
  consistent with local measurements~\cite{Riess2022SH0ES}.

  The correction \(\delta_{\mathrm{eff}}\) depends on the depth and spatial extent of
  the void, as well as on the characteristic saturation scale \(a_\star\).
  Denser or more weakly underdense regions remain close to the unsaturated regime
  and yield \(\delta_{\mathrm{eff}} \approx 0\).

  Importantly, this mechanism does not require any additional energy component
  or modification of the background Friedmann equations.
  The global expansion history remains unchanged.
  Only the local inference of the expansion rate is affected.

  Within this interpretation, the enhanced expansion inferred in cosmic voids
  reflects a reduced effective geometric constraint rather than an additional
  energy density.
  In regions where the density of contributing relational structure is low,
  the kinematic inference based on sparse mass distributions operates closer to its
  saturation limit.
  This induces a systematic environment-dependent bias in the locally inferred
  kinematic relations.
